% Shared content for the Final Exam.
% \input by BOTH final.tex (blank) and final-solutions.tex (answer key).
% Do NOT compile this file directly.
%
% The exam class hides \begin{solution}...\end{solution} unless the wrapper
% file calls \printanswers.

\begin{questions}

\question[10]
State and prove that a strictly concave function \(f:\R\to\R\) has at most one
maximizer.
\begin{solution}
Suppose \(x \ne y\) both maximize \(f\), so \(f(x) = f(y) = \max f\). Strict
concavity gives \(f\!\left(\tfrac{x+y}{2}\right) > \tfrac12 f(x) + \tfrac12 f(y)
= \max f\), a contradiction. Hence the maximizer is unique.
\end{solution}

\question[15]
A consumer has utility \(u(x_1,x_2) = x_1^{\alpha} x_2^{1-\alpha}\) with
\(\alpha \in (0,1)\), income \(m\), and prices \(p_1,p_2\).
\begin{parts}
  \part[6] Derive the Marshallian demands \(x_1(p,m)\) and \(x_2(p,m)\).
  \part[5] Compute the indirect utility function \(v(p,m)\).
  \part[4] Verify Roy's identity for good~1.
\end{parts}
\begin{solution}
Cobb--Douglas shares: \(x_1 = \dfrac{\alpha m}{p_1}\), \(x_2 = \dfrac{(1-\alpha)m}{p_2}\).
Then \(v(p,m) = \left(\tfrac{\alpha}{p_1}\right)^{\alpha}
\left(\tfrac{1-\alpha}{p_2}\right)^{1-\alpha} m\), and Roy's identity returns \(x_1\).
\end{solution}

\end{questions}
